% page 123 du fac-simile (image p-122 du scan)
% bloc 165.1 x 237.7 mm ; marge gauche 36.9 mm
\pgc{15.240mm}{0.000mm}{
\l{\cel{53}113}
\l{}
\l{}
\l{}
\l{\textquotedbl{}dis\textquotedbl{}. (\sou{Muziko)}\cel{1}. La noto duesma di la \textquotedbl{}do\textquotedbl{}-gamo, diezoza.}
\l{}
\l{dis-\cel{2}Prefixo qua indikas ideo di semeso hike ed ibe, hazarde.}
\l{}
\l{disenterio. (\sou{patol.}) Inflamuro di qua la fonto esas en la grosa}
\l{\cel{3}intestino, qua karakterizesas da ofta kaki di materio mukatra,}
\l{\cel{3}sangoza, akompanata da koliki forta e da doloriganta tenseso}
\l{\cel{3}di la regiono anusala qua bezonigas kako-probo (olqua esas}
\l{\cel{3}sen-sucesa). - DEFIRS.}
\l{}
\l{\sou{disertar}. (\sou{trans.pri}) Prizentar ideo-developo pri temo, pri}
\l{\cel{3}doktrino-punto. - DEFIRS.}
\l{}
\l{\sou{disho}. Nutrivo destinita a prizentesar lor repasto. - E.}
\l{}
\l{\sou{disidenta}. Qua ne konkordas kun la religiani maxim multa, pri}
\l{\cel{3}ula punto, dogmato. - DEFIRS.}
\l{}
\l{\sou{disimular}. (trans.)\cel{2}Celar sua pensi, opinioni, sentimenti, o}
\l{\cel{3}fingar altrai. - EFIS.}
\l{}
\l{\sou{disipar}. (\sou{trans}.) Ofte spensar abunde ed ecese de sua havajo, po}
\l{\cel{3}kozi frivola e sterila. (\sou{Metaf}.)Onu\cel{1}\sou{dis}ipas sua energio, sua}
\l{\cel{3}atenco ye multa objekti e skopi. - EFIS.}
\l{}
\l{\sou{disjuntiva}. I. (\sou{logiko}). Qua indikas la separeso di la du termini}
\l{\cel{3}di propoziciono. - II. (\sou{gram}.) Qua indikas ke la vorto, pri}
\l{\cel{3}lua senco, de la vorto qua sequas lu.}
\l{}
\l{\sou{disko}. I. (en\cel{1}\sou{la epoki antiqua}) Cirklajo pezoza, fera o petra,}
\l{\cel{3}quan onu probis lansar maxim adfore lor la exerci,la ludi}
\l{\cel{3}atletala. - II. Nomo di omna\cel{1}\sou{korp}o plata e cirkla (kom ex.:}
\l{\cel{3}disko fonografala, e c.) - DEFIRS.}
\l{}
\l{\sou{diskontar}. (\sou{trans.}) (\sou{Banko})\cel{2}Anticipe pagar ad ulu la pekunio-}
\l{\cel{3}quanto qua debesas a lu, prece di retenajo kalkulita segun}
\l{\cel{3}la interesti di la pekunio-quanto dum la tempo qua fluos til}
\l{\cel{3}la pago-dato normala, konvencionita (ye procento qua varias}
\l{\cel{3}segun la kurso di la merkato). - DEFIRS.}
\l{}
\l{\sou{disk}reta. Qua evitas jenar o tedar la kunhomi per sua demandi,}
\l{\cel{3}o nocar li per agi neprudenta (ex.: revelo di sekretaji).-}
\l{\cel{3}DEFIRS.}
\l{}
\l{\sou{diskriminanto}. (\sou{algebro}). I. Funciono di la koeficiento di}
\l{\cel{3}equaciono di la grado duesma, qua posibligas dicernar ka}
\l{\cel{3}la kurvo reprezentata da la equaciono en koordinati ortangula}
\l{\cel{3}esas kurvo vera, o sistemo de du rekti, - en la kazo di}
\l{\cel{3}elipso - kad ica esas reala od imaginala. - II. Dicesas anke}
\l{\cel{3}pri la equaciono B2 - E ac - a, b, c,\cel{2}esante la koeficienti}
\l{\cel{3}di ax2 + bx + c = 0.\cel{2}-\cel{2}DEF.}
}
